\section{Introduction}\label{sec:intro}

Tensor programs in PyTorch fail in two characteristic ways that
unit tests rarely catch: (i) the forward pass crashes deep inside
an autograd-traced subgraph because two operand shapes disagree
under a particular batch size, sequence length, or image
resolution; and (ii) the backward pass silently produces a wrong
gradient --- e.g.\ a tensor that should require a gradient does
not, because an in-place op or a \verb|.detach()| broke the tape
along one branch. Both classes manifest only at runtime, on a
specific device, often only on a specific input.

The v4 release of \textsc{TensorGuard} addressed (i) for a typed
AST fragment of \verb|nn.Module.forward| methods using symbolic
shape abstract interpretation with a Lean-checked operator core.
That formulation has two structural shortcomings. First, the
``shape'' analysis and the ``which tensors require grad'' analysis
were two ad-hoc passes over the same AST, with no shared
metatheory; the soundness statement covered only the shape pass.
Second, the analyzer treated each \verb|nn.Module| as a closed
program: a subclass that overrode \verb|forward| in a way the
parent's analysis had relied on broke the parent's verdict
silently, and there was no syntactic obligation marking the
breakage.

This paper presents a v5 rewrite around a single refinement-typed
calculus that subsumes both passes, makes module composition an
assume/guarantee discipline checkable at the class boundary, and
extends the verification target from forward shapes to
backward-pass gradient flags. Concretely, our contributions are:

\begin{enumerate}\itemsep0pt
\item \textbf{A refinement-typed shape-and-grad calculus}
  (\Cref{sec:calculus}). Tensor types are
  $\tau = \mathsf{Tensor}\{s : \mathsf{Shape},\ g :
  \mathsf{Bool}\mid \varphi\}$ where $s$ is a symbolic shape over
  a free commutative semiring on dim symbols, $g$ is a static
  gradient-required flag, and $\varphi$ is a Z3-decidable
  refinement constraining both. Operators have refinement-typed
  signatures whose preconditions are discharged by the same Z3
  procedure, eliminating the previous shape/grad pass split.
  Soundness is stated and proved against a small-step PyTorch
  semantics that models autograd tape construction explicitly.

\item \textbf{Assume/guarantee module composition}
  (\Cref{sec:assume-guarantee}). Each \verb|nn.Module| class is
  given a \emph{module contract}
  $\mathsf{A}\Rightarrow\mathsf{G}$ where $\mathsf{A}$
  refinement-types the inputs the class promises to accept and
  $\mathsf{G}$ refinement-types the outputs it promises to
  produce. Composition is verified locally: a caller need only
  show its callsite shape/grad refinements imply the callee's
  $\mathsf{A}$, and may then assume $\mathsf{G}$. Subclassing
  must refine the contract contravariantly on $\mathsf{A}$ and
  covariantly on $\mathsf{G}$; violations are a checkable
  syntactic obligation, not a silent re-analysis.

\item \textbf{An autograd-aware backward verifier}
  (\Cref{sec:impl-backward}). Using the grad-flag refinement, the
  analyzer constructs an abstract tape mirroring PyTorch's
  autograd graph and checks that, for every parameter the user
  marked \verb|requires_grad=True|, there exists a path of
  grad-carrying ops from that parameter to the loss. We catch the
  three canonical silent-zero-grad bugs (in-place op on a leaf,
  \verb|.detach()| in the loss path, \verb|torch.no_grad()|
  scoping error) at static time, with witnesses pointing at the
  edge that severed the tape.

\item \textbf{A Dynamo-guard correspondence and a benchmark
  scaled to that correspondence} (\Cref{sec:eval}). We prove that
  every refinement \textsc{TensorGuard} discharges at static time
  is implied by a guard that PyTorch's TorchDynamo would have
  installed at trace time on the same input
  (\Cref{thm:dynamo-corr}); conversely, every guard Dynamo
  installs whose body is in our supported fragment is a
  refinement we discharge. Empirically we report
  \textbf{(a)} an end-to-end benchmark of $\geq 150$ verification
  blocks drawn from torchvision, HuggingFace Transformers,
  timm, and a curated bug-injection suite, with per-block
  Dynamo-guard agreement rates;
  \textbf{(b)} an extended Lean parity development covering $20$
  operator transfer functions ($0$ \verb|sorry|, no
  \verb|mathlib|); and
  \textbf{(c)} a head-to-head comparison against
  \verb|torch.fx|+ShapeProp, \verb|FakeTensorMode|, and
  \verb|torch.export| under the v5 evaluation protocol.
\end{enumerate}

A reproducibility appendix (\Cref{app:repro}) lists every driver
script, JSON output, and Lean module backing the numbers in this
paper.
